% !TEX program = xelatex
% 공통수학2 · Ⅱ. 집합과 명제 · 1. 집합 · 01 집합 · 1/2차시
\documentclass[11pt,a4paper]{article}
\usepackage{kotex}
\usepackage{iftex}
\ifPDFTeX\else
  \setmainhangulfont{Noto Serif CJK KR}
  \setsanshangulfont{Noto Sans CJK KR}
\fi
\usepackage[margin=2.1cm]{geometry}
\usepackage{parskip}
\usepackage{enumitem}
\usepackage{array}
\usepackage{tabularx}
\usepackage{booktabs}
\usepackage{xcolor}
\usepackage{amssymb}
\usepackage{tikz}
\usepackage[most]{tcolorbox}
\usepackage{fancyhdr}
\usepackage{titlesec}

\definecolor{goldc}{HTML}{B07C22}
\definecolor{goldstrong}{HTML}{8A5F16}
\definecolor{indigoc}{HTML}{2B3B57}
\definecolor{indigosoft}{HTML}{6E82A6}
\definecolor{linec}{HTML}{D6DACB}
\definecolor{surface2}{HTML}{E4E8DC}
\definecolor{notebg}{HTML}{FBF3E3}
\definecolor{inksoft}{HTML}{52604F}

\titleformat{\section}{\normalfont\sffamily\bfseries\large\color{indigoc}}{\thesection}{0.6em}{}
\titlespacing{\section}{0pt}{16pt}{8pt}
\newtcolorbox{ideabox}{enhanced, breakable, colback=white, colframe=goldc, boxrule=0.5pt, leftrule=3pt, arc=2pt, left=10pt, right=10pt, top=8pt, bottom=8pt}
\newtcolorbox{calloutbox}{enhanced, breakable, colback=white, colframe=indigosoft, boxrule=0.5pt, leftrule=3pt, arc=2pt, left=10pt, right=10pt, top=7pt, bottom=7pt}
\newtcolorbox{notebox}{enhanced, breakable, colback=notebg, colframe=goldc, boxrule=0.5pt, leftrule=3pt, arc=2pt, left=10pt, right=10pt, top=7pt, bottom=7pt}
\newtcolorbox{answerbox}{enhanced, breakable, colback=white, colframe=linec, boxrule=0.5pt, arc=2pt, left=10pt, right=10pt, top=7pt, bottom=7pt}
\newcommand{\lessonbanner}[3]{%
  \begin{tcolorbox}[enhanced, colback=indigoc, colframe=indigoc, boxrule=0pt, arc=0pt,
    left=14pt, right=14pt, top=14pt, bottom=10pt, borderline south={2.4pt}{0pt}{goldc}, fontupper=\color{white}]
    {\sffamily\footnotesize\color{white!85!black} #1}\\[4pt]{\sffamily\bfseries\Large #2}\\[3pt]{\sffamily\small\color{white!88!black} #3}
  \end{tcolorbox}}
\pagestyle{fancy}\fancyhf{}\renewcommand{\headrulewidth}{0pt}
\fancyfoot[L]{\footnotesize\color{inksoft} 공통수학2 · 2026학년도 2학기 · 지평선고등학교}
\fancyfoot[R]{\footnotesize\color{inksoft} 교사용 지도안 -- 배포 금지}

\begin{document}
\lessonbanner{공통수학2 · Ⅱ. 집합과 명제 · 1. 집합}{01 집합 --- 1/2차시}{집합의 뜻과 표현 · 교사용 지도안 (2026학년도 2학기)}
\vspace{10pt}

\noindent\begin{tabularx}{\linewidth}{@{}>{\bfseries\color{indigoc}}p{2.6cm}X@{}}
\toprule
대상 & 고1 · 2개 반 · 반당 13명 \\
차시 & 50분 (도입 10 · 전개 30 · 정리 10) \\
성취기준 & 집합의 개념을 이해하고, 집합을 표현할 수 있다. \\
선행 학습 & 대단원 Ⅰ 전체 (도형의 방정식) \\
\bottomrule
\end{tabularx}

\section{학습 목표 \& Big Idea}
\begin{ideabox}
\textbf{\color{goldstrong}영속적 이해 (Big Idea)}\\
`분명한 기준'이 있어야 비로소 모임이 수학적 대상(집합)이 된다. 애매함을 걷어내고 명확한 판단 기준을 세우는 것이 수학적 사고의 출발점이다.
\vspace{4pt}
\small\color{inksoft} 핵심개념: \textbf{분명한 기준} \quad\textbullet\quad 핵심개념: \textbf{표현의 다양성} \quad\textbullet\quad 일반화: \textbf{같은 대상을 여러 방식으로 정확하게 표현할 수 있다}
\end{ideabox}
\begin{calloutbox}
\textbf{차시 학습 목표}
\begin{itemize}[leftmargin=1.4em, itemsep=2pt, topsep=4pt]
  \item 집합과 원소의 뜻을 이해하고, 집합인 것과 아닌 것을 구별할 수 있다.
  \item 집합을 원소나열법, 조건제시법, 벤 다이어그램으로 표현할 수 있다.
\end{itemize}
\end{calloutbox}

\section{질문의 연결}
\begin{tabularx}{\linewidth}{@{}>{\bfseries\color{indigoc}\small}p{2.4cm}X@{}}
사실적 질문 & 어떤 모임이 `집합'이 되려면 어떤 조건을 만족해야 할까? \\[6pt]
개념적 질문 & 같은 집합을 원소나열법과 조건제시법 두 가지로 나타낼 수 있는데, 언제 어느 표현이 더 유용할까? \\[6pt]
논쟁적 질문 & 재활용품 분류처럼 `기준이 분명한 모임'과 `귀여운 동물'처럼 `기준이 사람마다 다른 모임', 우리 사회의 규칙은 어느 쪽에 가까워야 할까? \\
\end{tabularx}

\section{수업 흐름}
\begin{tabularx}{\linewidth}{@{}>{\centering\arraybackslash}p{1.6cm}X>{\raggedright\arraybackslash\small}p{3.1cm}@{}}
\toprule
\textbf{단계} & \textbf{활동 \& 발문} & \textbf{ATL / VTR} \\
\midrule
\textcolor{goldstrong}{\textbf{도입}}\newline\footnotesize 10분 &
\textbf{마음공부 (2분)} --- 새 대단원 시작, 유무념 대조.\newline
\textbf{생각 열기 (8분)} --- 멸종 위기 야생 생물 목록(늑대, 수달, 혹고니, 수원청개구리, 감돌고기) 제시 $\to$ ``포유류에 속하는 동물을 모두 말해 보자''(기준 명확) vs ``귀여운 동물을 모두 말해 보자''(기준 불명확, 학생마다 답이 다름을 확인). &
ATL: 사고(기준의 명확성 비교)\newline VTR: See-Think-Wonder \\
\addlinespace
\textcolor{goldstrong}{\textbf{전개}}\newline\footnotesize 30분 &
\textbf{개념 형성 (10분)} --- 집합·원소 정의 $\to$ 문제 1(집합 여부 판별) $\to$ $\in$, $\notin$ 기호 $\to$ 문제 2.\newline
\textbf{표현법 (20분)} --- 원소나열법, 조건제시법을 각각 도입(8의 약수 예시) $\to$ 문제 3(상호 변환) $\to$ 벤 다이어그램 소개 $\to$ 문제 4(짝 활동). &
ATT: 촉진자 역할\newline ATL: 사고·표현 전환 \\
\addlinespace
\textcolor{goldstrong}{\textbf{정리}}\newline\footnotesize 10분 &
\textbf{개념 연결 질문 (5분)} --- ``세 가지 표현법(원소나열법, 조건제시법, 벤 다이어그램) 중 나는 언제 어느 것을 쓰고 싶은가?''\newline
\textbf{성찰 (5분)} --- 감사 일기 1문장 + 다음 차시 예고(원소의 개수, 유한집합·무한집합·공집합). &
마음공부 · 감사 일기 \\
\bottomrule
\end{tabularx}

\begin{calloutbox}
\textbf{지도상의 유의점} --- $a\in A$를 $A\in a$로 잘못 쓰지 않도록 순서에 주의시킨다. 원소나열법에서 원소를 나열하는 순서는 상관없지만, 같은 원소를 중복해서 쓰지 않는다는 점을 강조한다. 조건제시법 $\{x\mid \sim\}$에서 문자 $x$는 다른 문자로 바꿔도 됨을 언급한다.
\end{calloutbox}

\section{형성평가 \& 답안}
\begin{answerbox}
\textbf{문제 1} \quad 집합인 것을 찾고 원소 구하기\\
\small ⑴ 9의 약수 ⑵ 인기 있는 수영 선수 ⑶ 높은 빌딩 ⑷ 방정식 $4x-5=0$의 해\\[4pt]
\textbf{\color{goldstrong}답} \; 집합: ⑴, ⑷ (⑵, ⑶은 기준 불명확) / ⑴의 원소: $1,3,9$ / ⑷의 원소: $\dfrac54$
\end{answerbox}
\begin{answerbox}
\textbf{문제 2} \quad 실수 전체의 집합 $R$: ⑴ $0$ \quad ⑵ $\sqrt[3]{21}$ \quad ⑶ $-i$\\[4pt]
\textbf{\color{goldstrong}답} \; ⑴ $\in$ \quad ⑵ $\in$ \quad ⑶ $\notin$
\end{answerbox}
\begin{answerbox}
\textbf{문제 3} \quad 표현법 상호 변환\\
\small ⑴ $\{1,2,3,4,6,12\}$ ⑵ $\{5,10,15,\dots\}$ ⑶ $\{x\mid(x-3)(x-5)=0\}$ ⑷ $\{x\mid x$는 100 이하의 홀수$\}$\\[4pt]
\textbf{\color{goldstrong}답} \; ⑴ $\{x\mid x$는 12의 약수$\}$ \quad ⑵ $\{x\mid x$는 5의 배수$\}$ \quad ⑶ $\{3,5\}$ \quad ⑷ $\{1,3,5,\dots,99\}$
\end{answerbox}

\section{성취수준 (이 차시 범위)}
\begin{tabularx}{\linewidth}{@{}>{\bfseries\color{goldstrong}}p{0.9cm}X@{}}
\toprule
A & 집합의 개념을 설명하고, 집합을 다양한 방법으로 표현할 수 있다. \\
B & 집합의 개념을 이해하고, 집합을 다양한 방법으로 표현할 수 있다. \\
C & 집합의 개념을 알고, 집합을 표현할 수 있다. \\
D & 집합인 것과 아닌 것을 구분하고, 집합을 표현할 수 있다. \\
E & 집합인 것과 아닌 것을 구분할 수 있다. \\
\bottomrule
\end{tabularx}

\section{다음 차시}
\begin{calloutbox}
\textbf{01 집합 --- 2/2차시.} 유한집합, 무한집합, 공집합의 개념과 원소의 개수 $n(A)$를 다룬다.
\end{calloutbox}
\end{document}
